% BHU PYQ -- Political Science: Comparative Political Institutions-I (UK & Canada) (2025-26 NEP)
\documentclass[11pt,a4paper]{article}
\usepackage[a4paper,top=2.5cm,bottom=2.5cm,left=2.8cm,right=2.8cm,headheight=14pt]{geometry}
\usepackage{amsmath,amssymb}
\usepackage{fontspec}
\setmainfont{Kohinoor Devanagari}
\usepackage{microtype,fancyhdr,enumitem,hyperref,lastpage,parskip}

\hypersetup{colorlinks=true,urlcolor=black,linkcolor=black,pdftitle={POLMJ32: Comparative Political Institutions-I (UK & Canada) -- BHU NEP 2025-26}}

\pagestyle{fancy}\fancyhf{}
\renewcommand{\headrulewidth}{0.4pt}\renewcommand{\footrulewidth}{0.4pt}
\fancyhead[L]{\small \textit{Banaras Hindu University}}
\fancyhead[R]{\small \textit{POLMJ32: UK \& Canada Institutions (2025-26)}}
\fancyfoot[C]{\small Page \thepage\ of \pageref{LastPage}}

\newlist{parts}{enumerate}{2}
\setlist[parts,1]{label=(\alph*),leftmargin=*,itemsep=4pt,topsep=3pt}
\setlist[parts,2]{label=(\roman*),leftmargin=*,itemsep=2pt,topsep=2pt}

\newcommand{\pts}[1]{\hfill \textit{[#1]}}

\begin{document}

\begin{center}
  {\large\bfseries BANARAS HINDU UNIVERSITY}\\[2pt]
  {\normalsize B. A. (Hons.) Semester III Examination, 2025-26 (NEP)}\\[6pt]
  \rule{0.8\linewidth}{0.5pt}\\[6pt]
  {\large\bfseries POLITICAL SCIENCE}\\[4pt]
  {\large\bfseries Paper Code: POLMJ-32 : Comparative Political Institutions-I (UK \& Canada)}\\[4pt]
  {\normalsize Time: 2:30 Hours \hfill Full Marks: 70}\\[2pt]
  {\small REF NO. 2025-26/D-III/077}
\end{center}
\rule{\linewidth}{0.4pt}
\smallskip

\noindent \textit{Instructions: Attempt all questions of Sections A, B and C. Write your Roll No.\ at the top immediately on receipt of this question paper.}

\smallskip
\noindent \textit{सभी खण्डों अ, ब तथा स के प्रश्नों के उत्तर दीजिए।}

\bigskip

\section*{SECTION -- A / खण्ड -- अ}
\noindent \textit{Note: Short Answer Type Questions. Answer each question within 50 words.}\pts{5 $\times$ 2 = 10}

\noindent \textit{प्रत्येक प्रश्न का उत्तर 50 शब्दों के भीतर दीजिए। प्रत्येक प्रश्न 2 अंकों का है।}

\begin{enumerate}[label=\textbf{\arabic*.},leftmargin=*,itemsep=8pt]
  \item 
  \begin{parts}
    \item Rule of Law / विधि का शासन
    \item House of Lords / हाउस ऑफ लॉर्ड्स
    \item Canadian Federalism / कनाडाई संघवाद
    \item Shadow Cabinet / छाया मंत्रिमंडल
    \item Governor General of Canada / कनाडा के गवर्नर जनरल
  \end{parts}
\end{enumerate}

\bigskip

\section*{SECTION -- B / खण्ड -- ब}
\noindent \textit{Note: Answer the following questions in about 250 words each.}\pts{3 $\times$ 10 = 30}

\noindent \textit{निम्नलिखित प्रश्नों का उत्तर लगभग 250 शब्दों में दीजिए। प्रत्येक प्रश्न 10 अंकों का है।}

\begin{enumerate}[label=\textbf{\arabic*.},leftmargin=*,start=2,itemsep=12pt]

  \item Discuss the major conventions of the British Constitution.\pts{10}
  
  ब्रिटिश संविधान की प्रमुख अभिसमयों (परम्पराओं) की चर्चा कीजिए।

  \begin{center}\textbf{OR / अथवा}\end{center}

  Explain the powers and position of the Speaker of the House of Commons in the UK.\pts{10}
  
  ब्रिटेन के हाउस ऑफ कॉमन्स के अध्यक्ष (स्पीकर) की शक्तियों एवं स्थिति की व्याख्या कीजिए।

  \item Describe the structure of the political executive in Canada and the functions of the Prime Minister of Canada.\pts{10}
  
  कनाडा में राजनीतिक कार्यपालिका की संरचना तथा कनाडा के प्रधानमंत्री के कार्यों का वर्णन कीजिए।

  \begin{center}\textbf{OR / अथवा}\end{center}

  ``Rule of Law is a distinctive characteristic of the English Constitution.'' Comment.\pts{10}
  
  ``विधि का शासन ब्रिटिश संविधान का एक विशिष्ट लक्षण है।'' टिप्पणी कीजिए।

  \item Write a short note on the role of Prime Minister in the British political system.\pts{10}
  
  ब्रिटिश राजनीतिक व्यवस्था में प्रधानमंत्री की भूमिका पर एक संक्षिप्त टिप्पणी लिखिए।

  \begin{center}\textbf{OR / अथवा}\end{center}

  Examine the nature of the party system in Canada.\pts{10}
  
  कनाडा में दलीय व्यवस्था की प्रकृति का परीक्षण कीजिए।

\end{enumerate}

\bigskip

\section*{SECTION -- C / खण्ड -- स}
\noindent \textit{Note: Long Answer Type Questions. Answer all questions in detail.}\pts{2 $\times$ 15 = 30}

\noindent \textit{सभी प्रश्नों के उत्तर विस्तृत रूप से दीजिए। प्रत्येक प्रश्न 15 अंकों का है।}

\begin{enumerate}[label=\textbf{\arabic*.},leftmargin=*,start=5,itemsep=14pt]

  \item Critically examine Parliamentary Sovereignty in the United Kingdom.\pts{15}
  
  ब्रिटेन (यूके) में संसदीय संप्रभुता का समालोचनात्मक परीक्षण कीजिए।

  \begin{center}\textbf{OR / अथवा}\end{center}

  Analyze the nature and features of Canadian Federalism.\pts{15}
  
  कनाडाई संघवाद की प्रकृति एवं विशेषताओं का विश्लेषण कीजिए।

  \item Compare the powers and functions of the UK House of Commons and the Canadian House of Commons.\pts{15}
  
  ब्रिटिश हाउस ऑफ कॉमन्स तथा कनाडाई हाउस ऑफ कॉमन्स की शक्तियों एवं कार्यों की तुलना कीजिए।

  \begin{center}\textbf{OR / अथवा}\end{center}

  Discuss the powers, composition and relevance of the House of Lords in the British political system.\pts{15}
  
  ब्रिटिश राजनीतिक व्यवस्था में हाउस ऑफ लॉर्ड्स की शक्तियों, संरचना एवं प्रासंगिकता की विवेचना कीजिए।

\end{enumerate}

\end{document}
